%================================================================
% Poglavje 7: Sklep
%================================================================
\chapter{Sklep}
\label{ch:sklep}

V magistrskem delu smo predstavili in odprtokodno implementirali celovit cevovod za kardinalno omejeno optimizacijo portfelja, ki sodoben presečni napovedni model donosov (transformer tipa MASTER) povezuje s tremi klasičnimi metahevristikami, izvedenimi v jeziku C++. Ključna zasnovna odločitev je bila ločitev dveh virov težavnosti: oceno pričakovanih donosov in tveganja prepustimo napovednemu modelu, kombinatorično-nekonveksno optimizacijo pa metahevristikam, katerih kakovost neodvisno preverimo na standardnih primerih OR-Library. Tam se algoritmi približajo objavljeni neomejeni fronti na okoli odstotek pri manjših in nekaj odstotkov pri večjih univerzah, kar potrjuje, da je optimalnostna vrzel zanemarljiva.

Osrednje empirično vprašanje --- ali transformerski cevovod (model~$\vec{\mu}$ + model~$\vec{\Sigma}$) izboljša klasični zgodovinski cevovod --- smo ovrednotili v šestih tržnih režimih z združevanjem oken ($n\approx 254$) ter z dvema neposrednima primerjavama na protokolu objavljenih študij. Odgovor je pritrdilen, a niansiran. Transformer prekaša zgodovinsko referenco v petih od šestih režimov, je edini vir napovedi z zunaj-vzorčno pozitivno napovedno močjo in doseže najvišji tvegano prilagojeni donos; njegova prednost pred strojnimi in klasičnimi napovednimi modeli je statistično značilna po večih robustnih merilih, skupaj z ansamblom pa je edini, ki prestane na izbiro popravljeni Sharpov količnik. Na obeh tujih protokolih preseže objavljene reference. Poštena omejitev ostaja, da je prednost pred golim zgodovinskim povprečjem le ekonomsko smiselna, a statistično mejna; zato ansambel s postopkom Hedge predstavimo kot plast robustnosti, ki po zasnovi (meja regreta) v nobenem režimu ne zaostane bistveno za najboljšim virom, ne pa kot glavni rezultat.

\section{Nadaljnje delo}
\label{sec:future}

Delo je mogoče nadgraditi v več smereh. Napovedni del bi lahko obogatili z dodatnimi vhodnimi značilkami (temeljni in makroekonomski dejavniki) ter s pristopi, odpornejšimi na režimski premik --- npr.\ z instančno normalizacijo vhodov ali s cilji, ki neposredno optimizirajo tvegano prilagojeni donos. Kovariančno vejo bi lahko učili na realiziranih kovariancah iz visokofrekvenčnih podatkov. Na strani optimizacije bi bilo smiselno preizkusiti večciljne različice, ki bi celotno fronto trasirale v enem zagonu, ter dodatne realistične omejitve (transakcijski stroški neposredno v ciljni funkciji, sektorske omejitve). Za strožjo oceno prekomernega prilagajanja bi kazalo dodati kombinatorično-simetrično navzkrižno preverjanje. Nazadnje bi bilo cevovod vredno preizkusiti na več trgih, razredih sredstev in daljših obdobjih, da bi še utrdili oceno robustnosti napovednega signala skozi tržne režime.
